\documentclass[12pt]{article}
\usepackage[a3paper, margin=1in]{geometry}
\usepackage{amsmath, amssymb, amsthm, ulem}

\begin{document}

\section*{Domácí úkol 8.1}

Vypracoval: Daniel \textit{``Randál"} Ransdorf \hfill Podpis: \rule{4cm}{0.4pt}

\begin{enumerate}
  
  \item Obecnou matici, která splňuje podmínku jádra ze zadání, označme $A$. Označme si
    její sloupce $A = \left(\vec{a} \;\middle|\; \vec{b} \;\middle|\; \vec{c}\right)$.
    Nalezněme všechny matice $A$.
    Z definice jádra víme, že musí platit tento vztah:

    \begin{align*}
      &A \cdot \left(\begin{array}{c}
        1\\2\\3
      \end{array}\right) = \left(\begin{array}{c}
        0\\0\\0
      \end{array}\right) \\
      \Longleftrightarrow \quad
      &\vec{a} + 2\vec{b} + 3\vec{c} = \left(\begin{array}{c}
        0\\0\\0
      \end{array}\right) \\
      \Longleftrightarrow \quad
      &\vec{a} = -2\vec{b} - 3\vec{c} \\
      \Longleftrightarrow \quad
      &\left(\begin{array}{c}
        a_1 \\ a_2 \\ a_3
      \end{array}\right) = -2\left(\begin{array}{c}
        b_1 \\ b_2 \\ b_3
      \end{array}\right) - 3\left(\begin{array}{c}
        c_1 \\ c_2 \\ c_3
      \end{array}\right) \\
      \Longleftrightarrow \quad
      &\begin{cases}
        a_1 = -2b_1-3c_1 \\
        a_2 = -2b_2-3c_2 \\
        a_3 = -2b_3-3c_3 \\
      \end{cases} \quad | \;b_1, b_2, b_3, c_1, c_2, c_3 \in \mathbb{R}
    \end{align*}

    Dosaďme tedy parametricky vyjádřené hodnoty do matice $A$:
    
    \renewcommand{\arraystretch}{1.4}
    \begin{align*}
      A = \left(\vec{a} \;\middle|\; \vec{b} \;\middle|\; \vec{c}\right)
        = \left(\begin{array}{ccc}
          a_1 & b_1 & c_1 \\
          a_2 & b_2 & c_2 \\
          a_3 & b_3 & c_3
        \end{array}\right)
        = \left(\begin{array}{ccc}
          -2b_1-3c_1 & b_1 & c_1 \\
          -2b_2-3c_2 & b_2 & c_2 \\
          -2b_3-3c_3 & b_3 & c_3
        \end{array}\right) \quad | \;b_1, b_2, b_3, c_1, c_2, c_3 \in \mathbb{R}
    \end{align*}

    Rozložme matici na lineární kombinaci matic a tím získejme bázi:

    \begin{align*}
      A &= \left(\begin{array}{ccc}
          -2b_1-3c_1 & b_1 & c_1 \\
          -2b_2-3c_2 & b_2 & c_2 \\
          -2b_3-3c_3 & b_3 & c_3
        \end{array}\right) 
    \end{align*}
    \renewcommand{\arraystretch}{0.8}
    \small
    \begin{align*}
        &= \left(\begin{array}{ccc}
          -3c_1 & 0 & c_1 \\
          0 & 0 & 0 \\
          0 & 0 & 0
        \end{array}\right) + \left(\begin{array}{ccc}
          -2b_1 & b_1 & 0 \\
          0 & 0 & 0 \\
          0 & 0 & 0
        \end{array}\right) + \left(\begin{array}{ccc}
          0 & 0 & 0 \\
          -3c_2 & 0 & c_2 \\
          0 & 0 & 0
        \end{array}\right) + \left(\begin{array}{ccc}
          0 & 0 & 0 \\
          -2b_2 & b_2 & 0 \\
          0 & 0 & 0
        \end{array}\right) + \left(\begin{array}{ccc}
          0 & 0 & 0 \\
          0 & 0 & 0 \\
          -3c_3 & 0 & c_3 
        \end{array}\right) + \left(\begin{array}{ccc}
          0 & 0 & 0 \\
          0 & 0 & 0 \\
          -2b_3 & b_3 & 0 
        \end{array}\right) \\
        &= c_1\left(\begin{array}{ccc}
          -3 & 0 & 1 \\
          0 & 0 & 0 \\
          0 & 0 & 0
        \end{array}\right) + b_1\left(\begin{array}{ccc}
          -2 & 1 & 0 \\
          0 & 0 & 0 \\
          0 & 0 & 0
        \end{array}\right) + c_2\left(\begin{array}{ccc}
          0 & 0 & 0 \\
          -3 & 0 & 1 \\
          0 & 0 & 0
        \end{array}\right) + b_2\left(\begin{array}{ccc}
          0 & 0 & 0 \\
          -2& 1& 0 \\
          0 & 0 & 0
        \end{array}\right) + c_3\left(\begin{array}{ccc}
          0 & 0 & 0 \\
          0 & 0 & 0 \\
          -3 & 0 &  1
        \end{array}\right) + b_3\left(\begin{array}{ccc}
          0 & 0 & 0 \\
          0 & 0 & 0 \\
          -2 & 1 & 0 
        \end{array}\right) \quad | \;b_1, b_2, b_3, c_1, c_2, c_3 \in \mathbb{R}
    \end{align*}

    \normalsize

    Obecnou matice $A$ jsme rozložili na lineární kombinaci zjevně lineárně nezávislých maticí
    (každá obsahuje $1$ na místě, kde mají ostatní $0$, tím se žádná nedá vyjádřit LK ostatních).
    Tyto matice jsou také generátory, protože jimi dokážeme pomocí LK
    vyjádřit všechny matice splňující podmínku ze zadání.
    A tedy vytvořením posloupnosti z těchto matic zísáme bázi prostoru \textbf{W}.

    \[
      W = 
        Span \left(
          \left(\begin{array}{ccc}
          -3 & 0 & 1 \\
          0 & 0 & 0 \\
          0 & 0 & 0
        \end{array}\right) , \left(\begin{array}{ccc}
          -2 & 1 & 0 \\
          0 & 0 & 0 \\
          0 & 0 & 0
        \end{array}\right) , \left(\begin{array}{ccc}
          0 & 0 & 0 \\
          -3 & 0 & 1 \\
          0 & 0 & 0
        \end{array}\right) , \left(\begin{array}{ccc}
          0 & 0 & 0 \\
          -2& 1& 0 \\
          0 & 0 & 0
        \end{array}\right) , \left(\begin{array}{ccc}
          0 & 0 & 0 \\
          0 & 0 & 0 \\
          -3 & 0 &  1
        \end{array}\right) , \left(\begin{array}{ccc}
          0 & 0 & 0 \\
          0 & 0 & 0 \\
          -2 & 1 & 0 
        \end{array}\right)
        \right)
    \]
    

\end{enumerate}

\end{document}
